% GENERATED FILE. Edit canonical lessons, then run npm run generate:books.
% canonical-source-hash: fnv1a64:037761f176ab5608
% canonical-lessons: ES-C630-economia, ES-C630-economico, ES-C630-exportacion, ES-C630-inflacion, ES-C630-comercio

\chapter{La economía}
\label{ch:es-la-economia}
\hlchaptermodality{630}

\begin{chapteropening}
\textbf{By the end of this chapter:} \emph{I can follow simple news about the economy, inflation, exports and trade.}

Complete the last lesson of chapter 630: I can follow simple news about the economy, inflation, exports and trade.
\end{chapteropening}

\section[la economía]{la economía — the economy; economics}

\label{lesson:ES-C630-economia}



\begin{quote}
Before the new one: say the Spanish for the investment, then the Spanish for to finance.
\end{quote}

\subsection*{You'll want to know: la economía}
\textbf{la economía} — \textquotedblleft{}the economy; economics\textquotedblright{}.

\begin{quote}
La economía va mejor.
\end{quote}

\begin{quote}
\emph{The economy is doing better.}
\end{quote}

Also the subject at university.

\begin{cousinweb}
\emph{economía} is a learned borrowing from Latin \emph{oeconomia}, taken from Greek \emph{oikonomia} 'household management', built from \emph{oikos} 'house' and \emph{nomos} 'custom, rule' (EtymDB).
\end{cousinweb}

\begin{grammarlens}[title={What you've built}]
One more word for telling, explaining and describing at B1.
\end{grammarlens}

\subsection*{Your turn}
\begin{itemize}
\raggedright
  \item \emph{Say it:} \emph{la economía}
  \item \emph{Say it:} \emph{la economía}, once more
  \item \emph{Say it:} a sentence of your own with \emph{la economía}
  \item \emph{Recall:} say the Spanish for the investment, then the Spanish for to finance, then say \emph{la economía} again
\end{itemize}

\begin{tcolorbox}[breakable,colback=teal!4,colframe=teal!35!black,title={Before you move on}]
What does la economía mean? (\textquotedblleft{}The economy; economics\textquotedblright{}.) Say it once more.
\end{tcolorbox}
\section[económico]{económico — economic; cheap}

\label{lesson:ES-C630-economico}



\begin{quote}
Before the new one: say the Spanish for to finance, then the Spanish for the economy.
\end{quote}

\subsection*{You'll want to know: económico}
\textbf{económico} — \textquotedblleft{}economic; cheap\textquotedblright{}.

\begin{quote}
Es un hotel económico.
\end{quote}

\begin{quote}
\emph{It is a cheap hotel.}
\end{quote}

It means cheap as well as economic: \emph{una opción económica}.

\begin{cousinweb}
\emph{económico} is a learned borrowing from Latin \emph{oeconomicus}, taken from Greek \emph{oikonomikos} 'skilled in household management', from \emph{oikos} 'house' (EtymDB).
\end{cousinweb}

\begin{grammarlens}[title={What you've built}]
One more word for telling, explaining and describing at B1.
\end{grammarlens}

\subsection*{Your turn}
\begin{itemize}
\raggedright
  \item \emph{Say it:} \emph{económico}
  \item \emph{Say it:} \emph{económico}, once more
  \item \emph{Say it:} a sentence of your own with \emph{económico}
  \item \emph{Recall:} say the Spanish for to finance, then the Spanish for the economy, then say \emph{económico} again
\end{itemize}

\begin{tcolorbox}[breakable,colback=teal!4,colframe=teal!35!black,title={Before you move on}]
What does económico mean? (\textquotedblleft{}Economic; cheap\textquotedblright{}.) Say it once more.
\end{tcolorbox}
\section[la exportación]{la exportación — the export}

\label{lesson:ES-C630-exportacion}



\begin{quote}
Before the new one: say the Spanish for the economy, then the Spanish for economic.
\end{quote}

\subsection*{You'll want to know: la exportación}
\textbf{la exportación} — \textquotedblleft{}the export\textquotedblright{}.

\begin{quote}
La exportación de aceite crece.
\end{quote}

\begin{quote}
\emph{Olive oil exports are growing.}
\end{quote}

\emph{Exportar} — to export; \emph{la importación} — the import.

\begin{cousinweb}
\emph{exportación} is built on \emph{exportar}, a learned borrowing from Latin \emph{exportare} 'to carry out', from \emph{ex} and \emph{portare} 'to carry' (EtymDB).
\end{cousinweb}

\begin{grammarlens}[title={What you've built}]
One more word for telling, explaining and describing at B1.
\end{grammarlens}

\subsection*{Your turn}
\begin{itemize}
\raggedright
  \item \emph{Say it:} \emph{la exportación}
  \item \emph{Say it:} \emph{la exportación}, once more
  \item \emph{Say it:} a sentence of your own with \emph{la exportación}
  \item \emph{Recall:} say the Spanish for the economy, then the Spanish for economic, then say \emph{la exportación} again
\end{itemize}

\begin{tcolorbox}[breakable,colback=teal!4,colframe=teal!35!black,title={Before you move on}]
What does la exportación mean? (\textquotedblleft{}The export\textquotedblright{}.) Say it once more.
\end{tcolorbox}
\section[la inflación]{la inflación — the inflation}

\label{lesson:ES-C630-inflacion}



\begin{quote}
Before the new one: say the Spanish for economic, then the Spanish for the export.
\end{quote}

\subsection*{You'll want to know: la inflación}
\textbf{la inflación} — \textquotedblleft{}the inflation\textquotedblright{}.

\begin{quote}
La inflación está muy alta este año.
\end{quote}

\begin{quote}
\emph{Inflation is very high this year.}
\end{quote}

When inflation goes up, everything costs more.

\begin{cousinweb}
\emph{inflación} is a learned borrowing from Latin \emph{inflatio} 'a swelling', built on \emph{inflare} 'to blow into', from \emph{flare} 'to blow' (EtymDB).
\end{cousinweb}

\begin{grammarlens}[title={What you've built}]
One more word for telling, explaining and describing at B1.
\end{grammarlens}

\subsection*{Your turn}
\begin{itemize}
\raggedright
  \item \emph{Say it:} \emph{la inflación}
  \item \emph{Say it:} \emph{la inflación}, once more
  \item \emph{Say it:} a sentence of your own with \emph{la inflación}
  \item \emph{Recall:} say the Spanish for economic, then the Spanish for the export, then say \emph{la inflación} again
\end{itemize}

\begin{tcolorbox}[breakable,colback=teal!4,colframe=teal!35!black,title={Before you move on}]
What does la inflación mean? (\textquotedblleft{}The inflation\textquotedblright{}.) Say it once more.
\end{tcolorbox}
\section[el comercio]{el comercio — the trade, the commerce; the shop}

\label{lesson:ES-C630-comercio}



\begin{quote}
Before the new one: say the Spanish for the export, then the Spanish for the inflation.
\end{quote}

\subsection*{You'll want to know: el comercio}
\textbf{el comercio} — \textquotedblleft{}the trade, the commerce; the shop\textquotedblright{}.

\begin{quote}
El comercio con Europa crece.
\end{quote}

\begin{quote}
\emph{Trade with Europe is growing.}
\end{quote}

\emph{Un comercio} is also a shop: \emph{los comercios del barrio}.

\begin{cousinweb}
\emph{comercio} is a learned borrowing from Latin \emph{commercium} 'trade, dealings', built from \emph{con} and \emph{merx} 'goods' (EtymDB).
\end{cousinweb}

\begin{grammarlens}[title={What you've built}]
One more word for telling, explaining and describing at B1.
\end{grammarlens}

\subsection*{Your turn}
\begin{itemize}
\raggedright
  \item \emph{Say it:} \emph{el comercio}
  \item \emph{Say it:} \emph{el comercio}, once more
  \item \emph{Say it:} a sentence of your own with \emph{el comercio}
  \item \emph{Recall:} say the Spanish for the export, then the Spanish for the inflation, then say \emph{el comercio} again
\end{itemize}

\begin{tcolorbox}[breakable,colback=teal!4,colframe=teal!35!black,title={Before you move on}]
What does el comercio mean? (\textquotedblleft{}The trade, the commerce; the shop\textquotedblright{}.) Say it once more.
\end{tcolorbox}
